
\subsubsection{6.1 Key Finding: Inverted Gradient Dynamics}

The gradient competition hypothesis assumes that spurious features dominate because they ''win'' the competition for gradient signal. Our measurements reveal the opposite: minority samples---those requiring invariant feature learning---produce 5.$7\times$ higher gradient norms than spurious-aligned samples.

This inversion is consistent with what gradient magnitude indicates. High gradients signal ongoing learning---samples that remain difficult continue generating parameter updates. Low gradients indicate completed learning---samples that the model has already solved contribute little to further updates.

\subsubsection{6.2 Alternative Mechanism: Convergence Speed}

We propose that spurious dominance arises from convergence speed rather than gradient magnitude:

\begin{enumerate}
\item Simpler patterns occupy flatter loss landscape regions that are easier to navigate
\item Networks converge to spurious solutions faster than to invariant solutions
\item Early convergence locks in spurious representations before core features develop
\item High minority gradients cannot overcome established spurious features because the network's effective learning rate on spurious-aligned samples is already near zero
\end{enumerate}

This reframing connects to Sharpness-Aware Minimization. If spurious solutions occupy sharp minima, SAM's sharpness penalty would naturally disfavor spurious features.

\subsubsection{6.3 Implications for Robustness Methods}

Methods targeting gradient rebalancing may be addressing the wrong mechanism. Our results show hard samples already dominate gradient flow---the problem is not gradient magnitude but convergence timing.

Loss landscape interventions may be more effective. Methods that target loss landscape geometry (SAM, weight averaging) rather than sample-level gradients may address the actual mechanism of spurious dominance.

\subsubsection{6.4 Limitations}

\textbf{Single dataset.} We test only on Waterbirds. Spurious correlation dynamics may differ on CelebA or Colored MNIST.

\textbf{Synthetic region masks.} Our GradCAM analysis uses spatial heuristics (upper 60\% spurious, lower 40\% core) rather than ground-truth segmentation masks.

\textbf{Layer4 gradients only.} Earlier layers may show different dynamics.

\textbf{ResNet-50 architecture.} Vision Transformers may exhibit different spurious learning dynamics.

\textbf{Convergence speed hypothesis not directly tested.} We propose convergence speed as an alternative mechanism based on interpretation of the gradient inversion, but did not directly measure per-group loss trajectories or loss landscape curvature.

\subsubsection{6.5 What Was Not Tested}

The original hypothesis included predictions about learning rate and batch size effects on spurious feature timing (H-M2, H-M3). These experiments were blocked by the H-M1 gate failure. The causal chain from LR/batch size to gradient noise to spurious timing cannot be evaluated because the gradient competition mechanism itself was falsified.
